\documentclass[11pt,a4paper]{article}
\usepackage[margin=1in]{geometry}
\usepackage{amsmath,amssymb}
\usepackage{parskip}
\usepackage{hyperref}

\title{\textbf{Pre-Experiment Overview:\\ Electron Spin Resonance (ESR)}}
\author{IBRAHIM H.I. ABUSHAWISH \\ Istanbul University, Department of Physics}
\date{}

\begin{document}
\maketitle

\section*{Purpose of the Experiment}
According to the lab manual, the primary goal of this experiment is to observe the resonance absorption of high-frequency electromagnetic energy by unpaired electron spins in a DPPH (Diphenylpicrylhydrazyl) free-radical sample. By tuning the external magnetic field and measuring the resonance conditions on an oscilloscope, we determine the resonant magnetic field strength, $B_r$, and calculate the precise Land\'e-$g$ factor of the free electron.

\section*{Importance and Significance}
This experiment provides direct, spectroscopic confirmation of intrinsic electron spin and the Zeeman effect. In DPPH, each molecule contains a single unpaired valence electron whose orbital angular momentum is largely quenched by molecular bonding ($L \approx 0$), meaning its magnetic dipole moment arises almost purely from intrinsic electron spin, $S = 1/2$. When an external static magnetic field, $B$, is applied, this degenerate spin state splits into two discrete energy levels corresponding to $m_s = +1/2$ and $m_s = -1/2$, separated by an energy difference $\Delta E = g \mu_B B$.

In the laboratory apparatus, the DPPH sample sits inside a radiofrequency (RF) probe coil positioned symmetrically between a pair of Helmholtz coils. The Helmholtz coils supply both a DC bias field and an AC modulation field to sweep through the resonance point. When the RF photon energy $h \nu$ matches the Zeeman splitting $\Delta E$, the electrons absorb energy and flip spins. This resonance absorption is detected as a sharp peak on the oscilloscope screen. By recording the coil current at resonance and using the Helmholtz geometry formula $B = \mu_0 \left(\frac{4}{5}\right)^{3/2} \frac{n I}{R}$, we evaluate the resonant magnetic field and extract the Land\'e-$g$ factor, confirming it closely matches the theoretical free electron value $g \approx 2.0023$.

\end{document}
